\documentclass[11pt]{article}
\usepackage{amsmath,amssymb,amsthm,amsfonts}
\usepackage{geometry}
\geometry{margin=1in}
\usepackage{algorithm}
\usepackage{algpseudocode}

\newtheorem{theorem}{Theorem}
\newtheorem{lemma}{Lemma}
\newtheorem{definition}{Definition}
\newtheorem{assumption}{Assumption}
\newtheorem{proposition}{Proposition}
\newtheorem{corollary}{Corollary}

\title{Convergence Analysis of the Proximal Gradient Method}
\author{Optimization Theory Report}
\date{\today}

\begin{document}
\maketitle

%=====================================================================
\section*{Proximal Gradient Method (Forward--Backward Splitting)}

We analyze the composite optimization problem
\[
\min_{x \in \mathbb{R}^n} \; F(x) := f(x) + g(x),
\]
where $f$ is convex and $L$-smooth (differentiable with Lipschitz-continuous gradient) and $g$ is convex, proper, and lower semicontinuous (possibly nonsmooth). The algorithm under study is the \emph{proximal gradient method}:
\[
x_{k+1} = \operatorname{prox}_{\eta g}\bigl(x_k - \eta \nabla f(x_k)\bigr).
\]

%=====================================================================
\section*{Mathematical Formulation}

\textbf{Proximal operator.} For a proper convex lsc function $g$ and step size $\eta > 0$,
\[
\operatorname{prox}_{\eta g}(v) := \arg\min_{u \in \mathbb{R}^n} \left\{ g(u) + \frac{1}{2\eta}\|u - v\|^2 \right\}.
\]

\textbf{Update rule.} Given $x_k$:
\begin{align}
y_k &= x_k - \eta \nabla f(x_k) \qquad &\text{(forward / gradient step)}\\
x_{k+1} &= \operatorname{prox}_{\eta g}(y_k) \qquad &\text{(backward / proximal step)}.
\end{align}

\textbf{Hyperparameters.}
\begin{itemize}
\item Step size $\eta > 0$, chosen as $\eta \le 1/L$.
\item Initialization $x_0 \in \mathbb{R}^n$.
\item Termination tolerance $\epsilon > 0$.
\end{itemize}

\textbf{Gradient mapping.} It is convenient to define the \emph{gradient mapping}
\[
G_\eta(x) := \frac{1}{\eta}\bigl(x - \operatorname{prox}_{\eta g}(x - \eta \nabla f(x))\bigr),
\]
so that the update can be written compactly as $x_{k+1} = x_k - \eta\, G_\eta(x_k)$. When $g \equiv 0$, this reduces to ordinary gradient descent with $G_\eta(x) = \nabla f(x)$.

%=====================================================================
\section*{Pseudocode}

\begin{algorithm}[H]
\caption{Proximal Gradient Method}
\begin{algorithmic}[1]
\Require Smooth convex $f$, convex $g$, step size $\eta \in (0, 1/L]$, initial point $x_0$, tolerance $\epsilon$
\State $k \gets 0$
\Repeat
    \State $y_k \gets x_k - \eta \nabla f(x_k)$ \Comment{forward step}
    \State $x_{k+1} \gets \operatorname{prox}_{\eta g}(y_k)$ \Comment{backward step}
    \State $k \gets k+1$
\Until{$\|x_k - x_{k-1}\| \le \epsilon \eta$ \quad (i.e. $\|G_\eta(x_{k-1})\|\le \epsilon$)}
\State \Return $x_k$
\end{algorithmic}
\end{algorithm}

%=====================================================================
\section*{Dry Run Example}

We take $f(w) = (w-3)^2$, $g \equiv 0$ (so $\operatorname{prox}_{\eta g} = \mathrm{Id}$ and the method reduces to gradient descent), $w_0 = 0$, $\eta = 0.1$. Here $\nabla f(w) = 2(w-3)$.

\textbf{Iteration 1.}
\begin{align*}
\nabla f(w_0) &= 2(0-3) = -6,\\
y_0 &= w_0 - \eta\nabla f(w_0) = 0 - 0.1(-6) = 0.6,\\
w_1 &= \operatorname{prox}_{\eta g}(y_0) = 0.6,\\
f(w_1) &= (0.6-3)^2 = (-2.4)^2 = 5.76.
\end{align*}

\textbf{Iteration 2.}
\begin{align*}
\nabla f(w_1) &= 2(0.6-3) = -4.8,\\
y_1 &= 0.6 - 0.1(-4.8) = 0.6 + 0.48 = 1.08,\\
w_2 &= 1.08,\\
f(w_2) &= (1.08-3)^2 = (-1.92)^2 = 3.6864.
\end{align*}

\textbf{Iteration 3.}
\begin{align*}
\nabla f(w_2) &= 2(1.08-3) = -3.84,\\
y_2 &= 1.08 - 0.1(-3.84) = 1.08 + 0.384 = 1.464,\\
w_3 &= 1.464,\\
f(w_3) &= (1.464-3)^2 = (-1.536)^2 = 2.359296.
\end{align*}

The iterates $0 \to 0.6 \to 1.08 \to 1.464$ converge monotonically toward the minimizer $w^\ast = 3$, with objective values $9 \to 5.76 \to 3.6864 \to 2.3593$ decreasing as predicted by the theory below.

%=====================================================================
\section*{Assumptions}

\begin{assumption}[Convexity of $f$]\label{as:convf}
$f:\mathbb{R}^n\to\mathbb{R}$ is convex and differentiable:
\[
f(y) \ge f(x) + \langle \nabla f(x), y-x\rangle \quad \forall x,y.
\]
\end{assumption}

\begin{assumption}[$L$-smoothness of $f$]\label{as:smooth}
$\nabla f$ is $L$-Lipschitz:
\[
\|\nabla f(x) - \nabla f(y)\| \le L\|x-y\| \quad \forall x,y.
\]
This implies the descent lemma:
\[
f(y) \le f(x) + \langle \nabla f(x), y-x\rangle + \tfrac{L}{2}\|y-x\|^2.
\]
\end{assumption}

\begin{assumption}[Regularity of $g$]\label{as:g}
$g:\mathbb{R}^n\to\mathbb{R}\cup\{+\infty\}$ is proper, closed, and convex.
\end{assumption}

\begin{assumption}[Solvability]\label{as:sol}
The set of minimizers of $F = f+g$ is nonempty; let $x^\ast$ be a minimizer with $F^\ast = F(x^\ast)$.
\end{assumption}

\begin{assumption}[Step size]\label{as:step}
$0 < \eta \le 1/L$.
\end{assumption}

%=====================================================================
\section*{Main Convergence Theorem}

\begin{theorem}[Sublinear rate, convex case]\label{thm:main}
Under Assumptions~\ref{as:convf}--\ref{as:step}, the proximal gradient iterates satisfy
\[
F(x_K) - F^\ast \le \frac{\|x_0 - x^\ast\|^2}{2\eta K}, \qquad K = 1,2,\dots
\]
In particular, with $\eta = 1/L$,
\[
F(x_K) - F^\ast \le \frac{L\|x_0 - x^\ast\|^2}{2K} = O\!\left(\frac{1}{K}\right).
\]
\end{theorem}

\begin{theorem}[Linear rate, strongly convex case]\label{thm:strong}
If, in addition, $f$ is $\mu$-strongly convex ($\mu>0$), then with $\eta = 1/L$,
\[
\|x_{k+1} - x^\ast\|^2 \le \left(1 - \frac{\mu}{L}\right)\|x_k - x^\ast\|^2,
\]
hence $\|x_K - x^\ast\|^2 \le (1-\mu/L)^K \|x_0 - x^\ast\|^2$, a linear (geometric) rate.
\end{theorem}

%=====================================================================
\section*{Theoretical Properties}

\begin{itemize}
\item \textbf{Stability.} The proximal operator is firmly nonexpansive; combined with the contraction property of the gradient step, the forward-backward map is averaged, giving Fej\'er-monotone iterates: $\|x_{k+1}-x^\ast\|\le\|x_k-x^\ast\|$.
\item \textbf{Hyperparameter sensitivity.} Convergence requires $\eta\le 1/L$. Larger $\eta$ within this range gives faster constants; $\eta>2/L$ may diverge. Backtracking line search can estimate $L$ adaptively.
\item \textbf{Comparison to baselines.} Reduces to gradient descent when $g\equiv0$ and to the proximal point method when $f\equiv0$. The $O(1/K)$ rate matches gradient descent; accelerated variants (FISTA) achieve $O(1/K^2)$.
\end{itemize}

%=====================================================================
\section*{Preliminaries (Definitions and Lemmas)}

\begin{lemma}[Nonexpansiveness of the proximal operator]\label{lem:nonexp}
For proper closed convex $g$ and any $\eta>0$, the operator $\operatorname{prox}_{\eta g}$ is firmly nonexpansive, hence $1$-Lipschitz:
\[
\|\operatorname{prox}_{\eta g}(u) - \operatorname{prox}_{\eta g}(v)\| \le \|u-v\| \qquad \forall u,v.
\]
\end{lemma}
\begin{proof}
Let $p = \operatorname{prox}_{\eta g}(u)$ and $q = \operatorname{prox}_{\eta g}(v)$.
\textbf{[Step 1]} By the optimality (first-order) condition of the prox-minimization, the subgradient inclusion gives
\[
u - p \in \eta\,\partial g(p), \qquad v - q \in \eta\,\partial g(q).
\]
\textbf{[Step 2]} By monotonicity of the subdifferential of the convex function $g$,
\[
\langle (u-p) - (v-q),\; p - q\rangle \ge 0.
\]
\textbf{[Step 3]} Rearranging,
\[
\langle u-v,\; p-q\rangle \ge \|p-q\|^2.
\]
\textbf{[Step 4]} By Cauchy--Schwarz, $\langle u-v, p-q\rangle \le \|u-v\|\,\|p-q\|$, so
\[
\|p-q\|^2 \le \|u-v\|\,\|p-q\| \implies \|p-q\| \le \|u-v\|. \qedhere
\]
\end{proof}

\begin{lemma}[Prox optimality / subgradient identity]\label{lem:proxopt}
With $x_{k+1} = \operatorname{prox}_{\eta g}(x_k - \eta\nabla f(x_k))$, define $G_k := G_\eta(x_k) = \tfrac{1}{\eta}(x_k - x_{k+1})$. Then
\[
G_k - \nabla f(x_k) \in \partial g(x_{k+1}).
\]
\end{lemma}
\begin{proof}
\textbf{[Step 5]} The optimality condition of $x_{k+1}=\arg\min_u \{g(u) + \tfrac{1}{2\eta}\|u - (x_k-\eta\nabla f(x_k))\|^2\}$ reads
\[
0 \in \partial g(x_{k+1}) + \tfrac{1}{\eta}\bigl(x_{k+1} - x_k + \eta\nabla f(x_k)\bigr).
\]
\textbf{[Step 6]} Hence
\[
\tfrac{1}{\eta}(x_k - x_{k+1}) - \nabla f(x_k) \in \partial g(x_{k+1}),
\]
i.e. $G_k - \nabla f(x_k) \in \partial g(x_{k+1})$. \qedhere
\end{proof}

\begin{lemma}[Descent lemma for the gradient mapping]\label{lem:descent}
Under Assumptions~\ref{as:smooth} and \ref{as:step} (i.e. $\eta\le1/L$), for any $z\in\mathbb{R}^n$,
\[
F(x_{k+1}) \le F(z) + \langle G_k, x_k - z\rangle - \frac{\eta}{2}\|G_k\|^2.
\]
\end{lemma}
\begin{proof}
\textbf{[Step 7]} Apply the descent lemma (Assumption~\ref{as:smooth}) at $x_{k+1} = x_k - \eta G_k$:
\[
f(x_{k+1}) \le f(x_k) + \langle \nabla f(x_k), x_{k+1}-x_k\rangle + \tfrac{L}{2}\|x_{k+1}-x_k\|^2.
\]
\textbf{[Step 8]} Substitute $x_{k+1}-x_k = -\eta G_k$ and use $\tfrac{L}{2}\eta^2 \le \tfrac{\eta}{2}$ (since $\eta\le1/L \Rightarrow L\eta\le1$):
\[
f(x_{k+1}) \le f(x_k) - \eta\langle \nabla f(x_k), G_k\rangle + \tfrac{\eta}{2}\|G_k\|^2.
\]
\textbf{[Step 9]} Convexity of $f$ (Assumption~\ref{as:convf}): for any $z$,
\[
f(x_k) \le f(z) + \langle \nabla f(x_k), x_k - z\rangle.
\]
\textbf{[Step 10]} Convexity of $g$ and Lemma~\ref{lem:proxopt}: since $G_k-\nabla f(x_k)\in\partial g(x_{k+1})$,
\[
g(x_{k+1}) \le g(z) - \langle G_k - \nabla f(x_k),\, z - x_{k+1}\rangle
= g(z) + \langle G_k - \nabla f(x_k),\, x_{k+1} - z\rangle.
\]
\textbf{[Step 11]} Add the three inequalities (Steps 8, 9, 10). First combine Steps 8 and 9:
\[
f(x_{k+1}) \le f(z) + \langle \nabla f(x_k), x_k - z\rangle - \eta\langle\nabla f(x_k),G_k\rangle + \tfrac{\eta}{2}\|G_k\|^2.
\]
\textbf{[Step 12]} Now add Step 10 and write $x_{k+1} - z = (x_k - z) - \eta G_k$:
\begin{align*}
F(x_{k+1}) &\le f(z)+g(z) + \langle \nabla f(x_k), x_k - z\rangle - \eta\langle\nabla f(x_k),G_k\rangle + \tfrac{\eta}{2}\|G_k\|^2 \\
&\quad + \langle G_k - \nabla f(x_k),\, (x_k-z) - \eta G_k\rangle.
\end{align*}
\textbf{[Step 13]} Expand the last inner product:
\[
\langle G_k - \nabla f(x_k), x_k - z\rangle - \eta\langle G_k-\nabla f(x_k), G_k\rangle.
\]
Collect the terms involving $\nabla f(x_k)$ and $G_k$:
\begin{itemize}
\item Coefficient of $\langle \nabla f(x_k), x_k-z\rangle$: $+1 - 1 = 0$ (cancels).
\item Coefficient of $\langle G_k, x_k - z\rangle$: $+1$.
\item Terms in $\langle\nabla f(x_k),G_k\rangle$: $-\eta + \eta = 0$ (cancels).
\item Terms in $\|G_k\|^2$: $\tfrac{\eta}{2} - \eta = -\tfrac{\eta}{2}$.
\end{itemize}
\textbf{[Step 14]} Therefore
\[
F(x_{k+1}) \le F(z) + \langle G_k, x_k - z\rangle - \frac{\eta}{2}\|G_k\|^2. \qedhere
\]
\end{proof}

%=====================================================================
\section*{Main Proof (Step-by-Step Derivation)}

\begin{proof}[Proof of Theorem~\ref{thm:main}]
\textbf{[Step 15]} Apply Lemma~\ref{lem:descent} with $z = x_k$:
\[
F(x_{k+1}) \le F(x_k) - \frac{\eta}{2}\|G_k\|^2 \le F(x_k),
\]
so the objective is monotonically nonincreasing.

\textbf{[Step 16]} Apply Lemma~\ref{lem:descent} with $z = x^\ast$:
\[
F(x_{k+1}) - F^\ast \le \langle G_k, x_k - x^\ast\rangle - \frac{\eta}{2}\|G_k\|^2.
\]

\textbf{[Step 17]} Complete the square. Using $x_{k+1} = x_k - \eta G_k$,
\begin{align*}
\|x_{k+1}-x^\ast\|^2 &= \|x_k - \eta G_k - x^\ast\|^2\\
&= \|x_k - x^\ast\|^2 - 2\eta\langle G_k, x_k - x^\ast\rangle + \eta^2\|G_k\|^2.
\end{align*}
Hence
\[
\langle G_k, x_k-x^\ast\rangle - \frac{\eta}{2}\|G_k\|^2
= \frac{1}{2\eta}\Bigl(\|x_k - x^\ast\|^2 - \|x_{k+1}-x^\ast\|^2\Bigr).
\]
\textbf{[Step 18]} (Verification.) Indeed,
\[
\frac{1}{2\eta}\Bigl(\|x_k-x^\ast\|^2 - \|x_{k+1}-x^\ast\|^2\Bigr)
= \frac{1}{2\eta}\Bigl(2\eta\langle G_k, x_k-x^\ast\rangle - \eta^2\|G_k\|^2\Bigr)
= \langle G_k, x_k-x^\ast\rangle - \frac{\eta}{2}\|G_k\|^2.
\]

\textbf{[Step 19]} Combining Steps 16 and 17,
\[
F(x_{k+1}) - F^\ast \le \frac{1}{2\eta}\Bigl(\|x_k - x^\ast\|^2 - \|x_{k+1}-x^\ast\|^2\Bigr).
\]

\textbf{[Step 20]} Sum over $k = 0,1,\dots,K-1$ (telescoping the right side):
\[
\sum_{k=0}^{K-1}\bigl(F(x_{k+1}) - F^\ast\bigr)
\le \frac{1}{2\eta}\Bigl(\|x_0 - x^\ast\|^2 - \|x_K - x^\ast\|^2\Bigr)
\le \frac{\|x_0 - x^\ast\|^2}{2\eta}.
\]

\textbf{[Step 21]} By Step 15, the sequence $\{F(x_{k+1})\}$ is nonincreasing, so $F(x_K) \le F(x_{k+1})$ for all $k \le K-1$. Therefore
\[
K\bigl(F(x_K) - F^\ast\bigr) \le \sum_{k=0}^{K-1}\bigl(F(x_{k+1}) - F^\ast\bigr) \le \frac{\|x_0 - x^\ast\|^2}{2\eta}.
\]

\textbf{[Step 22]} Dividing by $K$,
\[
F(x_K) - F^\ast \le \frac{\|x_0 - x^\ast\|^2}{2\eta K}.
\]
With $\eta = 1/L$ this gives $F(x_K)-F^\ast \le \dfrac{L\|x_0-x^\ast\|^2}{2K}$. \qedhere
\end{proof}

%=====================================================================
\section*{Rate Derivation (With Constants)}

From Step 22, to guarantee $F(x_K) - F^\ast \le \epsilon$ it suffices that
\[
\frac{\|x_0-x^\ast\|^2}{2\eta K} \le \epsilon
\quad\Longleftrightarrow\quad
K \ge \frac{\|x_0 - x^\ast\|^2}{2\eta\epsilon}.
\]
With $\eta = 1/L$, the iteration complexity is
\[
K = \left\lceil \frac{L\|x_0 - x^\ast\|^2}{2\epsilon}\right\rceil = O\!\left(\frac{L\|x_0-x^\ast\|^2}{\epsilon}\right).
\]
The explicit constant is $C = \tfrac{1}{2}L\|x_0-x^\ast\|^2$, so $F(x_K)-F^\ast \le C/K$.

%=====================================================================
\section*{Special Cases}

\paragraph{Case 1: $g\equiv 0$ (Gradient Descent).}
Then $\operatorname{prox}_{\eta g}=\mathrm{Id}$ and $G_k=\nabla f(x_k)$. Theorem~\ref{thm:main} recovers the classical gradient descent bound
\[
f(x_K) - f^\ast \le \frac{\|x_0-x^\ast\|^2}{2\eta K} = \frac{L\|x_0-x^\ast\|^2}{2K}.
\]

\paragraph{Case 2: $f\equiv 0$ (Proximal Point Method).}
Then $\nabla f=0$, $G_k=\tfrac1\eta(x_k-x_{k+1})$, and the update is $x_{k+1}=\operatorname{prox}_{\eta g}(x_k)$. The smoothness constant $L=0$ allows any $\eta>0$, and the same telescoping gives $g(x_K)-g^\ast\le \|x_0-x^\ast\|^2/(2\eta K)$.

\paragraph{Case 3: Strongly convex $f$ (Linear rate).}
\begin{proof}[Proof sketch of Theorem~\ref{thm:strong}]
\textbf{[Step 23]} With $\mu$-strong convexity, Step 9 strengthens to
\[
f(x_k) \le f(z) + \langle\nabla f(x_k), x_k - z\rangle - \frac{\mu}{2}\|x_k - z\|^2.
\]
\textbf{[Step 24]} Repeating Steps 10--14 with $z=x^\ast$ yields the refined bound
\[
F(x_{k+1}) - F^\ast \le \frac{1}{2\eta}\Bigl(\|x_k-x^\ast\|^2 - \|x_{k+1}-x^\ast\|^2\Bigr) - \frac{\mu}{2}\|x_k - x^\ast\|^2.
\]
\textbf{[Step 25]} Since $F(x_{k+1})-F^\ast\ge0$, dropping it and rearranging gives
\[
\|x_{k+1}-x^\ast\|^2 \le (1-\eta\mu)\|x_k-x^\ast\|^2.
\]
With $\eta=1/L$, $\;\|x_{k+1}-x^\ast\|^2 \le (1-\mu/L)\|x_k-x^\ast\|^2$, proving the geometric contraction. \qedhere
\end{proof}

%=====================================================================
\section*{Discussion}

The proof crucially used Lemma~\ref{lem:nonexp} (nonexpansiveness of $\operatorname{prox}$) implicitly through the optimality condition of Lemma~\ref{lem:proxopt}, which is equivalent to firm nonexpansiveness. The resulting $O(1/K)$ rate matches gradient descent and is optimal for first-order methods on smooth convex composite problems without acceleration; Nesterov-type momentum (FISTA) improves this to $O(1/K^2)$. The strongly convex case yields linear convergence with rate $(1-\mu/L)$, governed by the condition number $\kappa = L/\mu$.

\end{document}